\chapter{Monotone networks: universality for monotone functions}

\section{The threshold-network model (Mikulincer--Reichman)}

\begin{definition}[Heaviside gate]\label{def:heaviside}
  \lean{UniversalApproximation.Monotone.heaviside}\leanok
  $\mathrm{heaviside}(z)=1$ if $z\ge 0$ and $0$ otherwise.
\end{definition}

\begin{definition}[Activation stack]\label{def:actstack}
  \lean{UniversalApproximation.Monotone.ActStack}\leanok
  \uses{def:heaviside}
  An \texttt{ActStack} is an inductively-built chain of affine layers, each followed by a
  one-dimensional activation.
\end{definition}

\begin{definition}[Monotone network]\label{def:mononet}
  \lean{UniversalApproximation.Monotone.MonoNet}\leanok
  \uses{def:actstack}
  A \texttt{MonoNet} is an activation stack together with an affine read-out (weights
  \texttt{readW} and bias \texttt{readBias}).
\end{definition}

\begin{definition}[Domination gadget]\label{def:domination}
  \lean{UniversalApproximation.Monotone.dominationStack}\leanok
  \uses{def:actstack, def:heaviside}
  A two-layer threshold gadget whose output is the indicator $\indicator(x\ge x_j)$ of
  coordinatewise domination of a fixed point $x_j$.
\end{definition}

\begin{theorem}[Monotone interpolation, M--R Result 1]\label{thm:mono-interp}
  \lean{UniversalApproximation.Monotone.monotone_interpolation}\leanok
  \uses{def:mononet, def:heaviside, def:domination}
  For injective points $x$ and targets $y$ monotone along the coordinatewise order, some monotone
  depth-4 \texttt{MonoNet} agrees with $y$ on every $x_i$ exactly.
\end{theorem}
\begin{verbatim}
theorem monotone_interpolation (x : Fin n → (Fin d → ℝ)) (y : Fin n → ℝ)
    (hmono : ∀ i j, x i ≤ x j → y i ≤ y j) (hinj : Function.Injective x) :
    ∃ N : MonoNet d, N.IsMonotone ∧ N.depth = 4 ∧ ∀ i, N.toFun (x i) = y i
\end{verbatim}
\begin{proof}\leanok
  \uses{def:domination}
  The hidden layers compute, for each index $j$, the domination indicator
  $\indicator(x\ge x_j)$ via the two-layer gadget (\cref{def:domination}). Monotonicity of $y$
  along the order lets the read-out take a non-negative combination of these indicators that
  reproduces $y$ exactly on the injective points. The exact ($\varepsilon=0$) engine gives level
  sets on the nose, so the read-out is exact and the whole net is monotone of depth 4.
\end{proof}

\begin{theorem}[Monotone approximation, M--R Result 1]\label{thm:mono-approx}
  \lean{UniversalApproximation.Monotone.monotone_approximation}\leanok
  \uses{thm:mono-interp}
  Every continuous monotone $f$ on the unit cube is uniformly $\varepsilon$-approximated by a
  monotone depth-4 \texttt{MonoNet}.
\end{theorem}
\begin{verbatim}
theorem monotone_approximation {d : ℕ} (f : (Fin d → ℝ) → ℝ)
    (hf : ContinuousOn f (Set.Icc 0 1))
    (hmono : ∀ ⦃a b⦄, a ∈ Set.Icc (0 : Fin d → ℝ) 1 →
      b ∈ Set.Icc (0 : Fin d → ℝ) 1 → a ≤ b → f a ≤ f b)
    {ε : ℝ} (hε : 0 < ε) :
    ∃ N : MonoNet d, N.IsMonotone ∧ N.depth = 4 ∧
      ∀ x ∈ Set.Icc (0 : Fin d → ℝ) 1, |N.toFun x - f x| ≤ ε
\end{verbatim}
\begin{proof}\leanok
  \uses{thm:mono-interp}
  Sample $f$ on a fine grid of the cube and interpolate the grid values exactly with
  \cref{thm:mono-interp}. Uniform continuity of $f$ bounds the gap between grid nodes by
  $\varepsilon$; the interpolant is monotone of depth 4.
\end{proof}

\section{Beyond thresholds: saturating activations (Sartor et al.\ 2025)}

\begin{definition}[One-sided saturation, Def 3.3]\label{def:saturating}
  \lean{UniversalApproximation.Monotone.RightSaturating}
  \lean{UniversalApproximation.Monotone.LeftSaturating}\leanok
  $\sigma$ is \emph{right-saturating} ($\mathcal{S}^+$) if it has a finite limit at $+\infty$, and
  \emph{left-saturating} ($\mathcal{S}^-$) if it has a finite limit at $-\infty$. This admits
  unbounded activations such as ReLU on one side.
\end{definition}

\begin{definition}[Point reflection, Prop 3.8]\label{def:reflect}
  \lean{UniversalApproximation.Monotone.reflect}\leanok
  $\mathrm{reflect}(\sigma)(x) = -\sigma(-x)$. It is an involution, preserves monotonicity, and
  swaps the two saturation sides.
\end{definition}

\begin{lemma}[Saturation gain estimate, Lemma 3.6]\label{lem:sat-scaled}
  \lean{UniversalApproximation.Monotone.rightSaturating_scaled_approx}\leanok
  \uses{def:saturating}
  For a right-saturating $\sigma$ with limit $L$, and any $\varepsilon,m>0$, there is a gain
  $\Lambda$ so that for all $\lambda\ge\Lambda$ and $t\ge m$, $|\sigma(\lambda t)-L|\le\varepsilon$.
\end{lemma}
\begin{proof}\leanok
  Immediate from the definition of the limit at $+\infty$: for $t\ge m$, $\lambda t\to+\infty$ as
  $\lambda\to\infty$, so $\sigma(\lambda t)$ enters any $\varepsilon$-ball of $L$.
\end{proof}

\begin{lemma}[Interior continuity, approx value]\label{lem:approx-interior}
  \lean{UniversalApproximation.Monotone.approx_interior_value}\leanok
  If $\sigma$ is continuous at $b$, then for any $\varepsilon>0$ there is $\delta>0$ with
  $|\sigma(t)-\sigma(b)|\le\varepsilon$ whenever $|t-b|\le\delta$.
\end{lemma}
\begin{proof}\leanok
  A restatement of continuity at $b$ as an explicit $\varepsilon$--$\delta$ bound.
\end{proof}

\begin{proposition}[Weight-sign / saturation, Prop 3.10]\label{prop:sign-sat}
  \lean{UniversalApproximation.Monotone.prop_3_10_two_layer}\leanok
  \uses{def:reflect}
  Negating a layer's input and using the reflected activation $\mathrm{reflect}(\sigma)$ equals
  negating the next layer's weights and using $\sigma$: the weight-sign flip is absorbed by point
  reflection.
\end{proposition}
\begin{proof}\leanok
  \uses{def:reflect}
  Expand both layer compositions and apply $\mathrm{reflect}(\sigma)(x)=-\sigma(-x)$; the two
  sign flips cancel, matching the negated-weight layer applied to $\sigma$.
\end{proof}

\begin{theorem}[Saturating interpolation, Sartor Thm 3.5]\label{thm:sat-interp}
  \lean{UniversalApproximation.Monotone.saturating_interpolation}\leanok
  \uses{def:mononet, def:saturating, lem:sat-scaled, lem:approx-interior}
  For alternating monotone one-sided-saturating activations
  $[\sigma_1,\sigma_2,\sigma_3]=[\mathcal{S}^-,\mathcal{S}^+,\mathcal{S}^-]$, each non-constant,
  a monotone depth-4 \texttt{MonoNet} interpolates any monotone dataset to accuracy
  $\varepsilon$.
\end{theorem}
\begin{verbatim}
theorem saturating_interpolation {d n : ℕ} (x : Fin n → (Fin d → ℝ))
    (y : Fin n → ℝ) (hmono : ∀ i j, x i ≤ x j → y i ≤ y j)
    (hinj : Function.Injective x) (σ₁ σ₂ σ₃ : ℝ → ℝ)
    (hm₁ : Monotone σ₁) (hm₂ : Monotone σ₂) (hm₃ : Monotone σ₃)
    (hs₁ : LeftSaturating σ₁) (hs₂ : RightSaturating σ₂) (hs₃ : LeftSaturating σ₃)
    (hnc₁ : ∃ a b, σ₁ a < σ₁ b) (hnc₂ : ∃ a b, σ₂ a < σ₂ b)
    (hnc₃ : ∃ a b, σ₃ a < σ₃ b) {ε : ℝ} (hε : 0 < ε) :
    ∃ N : MonoNet d, N.IsMonotone ∧ N.depth = 4 ∧
      N.stack.activations = [σ₁, σ₂, σ₃] ∧ ∀ i, |N.toFun (x i) - y i| ≤ ε
\end{verbatim}
\begin{proof}\leanok
  \uses{thm:mono-interp, lem:sat-scaled, lem:approx-interior}
  Start from the threshold interpolant (\cref{thm:mono-interp}) and replace each Heaviside gate by
  a scaled saturating activation. With alternating sides $[\mathcal{S}^-,\mathcal{S}^+,
  \mathcal{S}^-]$, the gain estimate (\cref{lem:sat-scaled}) makes each gate approximate its
  threshold arbitrarily well, and interior continuity (\cref{lem:approx-interior}) controls values
  near the boundary; choosing the gain large enough yields $\varepsilon$-accuracy.
  \emph{Deviation note.} The paper's exact identity $g(x_i)=f(x_i)$ is the $\lambda\to\infty$
  idealization; the faithful finite-net statement is $\varepsilon$-approximate, and the
  non-degeneracy hypotheses are exactly $\sigma_k$ non-constant ($\exists a\,b,\ \sigma_k a <
  \sigma_k b$).
\end{proof}

\begin{proposition}[Non-positive weights, Prop 3.11]\label{prop:nonpos}
  \lean{UniversalApproximation.Monotone.nonpos_weight_universal}\leanok
  \uses{thm:sat-interp, def:reflect, prop:sign-sat}
  For a monotone, non-constant, right-saturating $\sigma$, any monotone dataset is
  $\varepsilon$-interpolated by a depth-4 network with all non-positive weights and the single
  activation $\sigma$ on its three hidden layers --- whose denotation is nonetheless monotone.
\end{proposition}
\begin{verbatim}
theorem nonpos_weight_universal {d n : ℕ} (x : Fin n → (Fin d → ℝ))
    (y : Fin n → ℝ) (hmono : ∀ i j, x i ≤ x j → y i ≤ y j)
    (hinj : Function.Injective x) (σ : ℝ → ℝ) (hmσ : Monotone σ)
    (hsat : RightSaturating σ) (hnc : ∃ a b, σ a < σ b)
    {ε : ℝ} (hε : 0 < ε) :
    ∃ N : MonoNet d, N.stack.WeightsNonpos ∧ (∀ i, N.readW i ≤ 0) ∧
      N.stack.activations = [σ, σ, σ] ∧ N.depth = 4 ∧ Monotone N.toFun ∧
      ∀ i, |N.toFun (x i) - y i| ≤ ε
\end{verbatim}
\begin{proof}\leanok
  \uses{thm:sat-interp, prop:sign-sat}
  Apply \cref{thm:sat-interp} with activations $[\mathrm{reflect}\,\sigma,\sigma,
  \mathrm{reflect}\,\sigma]=[\mathcal{S}^-,\mathcal{S}^+,\mathcal{S}^-]$, then use the Prop 3.10
  sign-flip identities (\cref{prop:sign-sat}) to turn every layer into $\sigma$ with negated
  weights; the final output sign flip is absorbed into the read-out, so the denotation is
  unchanged and stays monotone.
\end{proof}
